\section{Level 4: Counterfactual Outcome Reasoning}
\label{sec:level4}

\subsection{Capability Definition and Adjacent Boundary}

Level~4 asks what would have happened to the same patient or state under an action other than the one observed. Formally, the target is a potential outcome or trajectory such as $P(Y_a\mid X,A{=}a')$, not merely a sampled future labeled with action $a$. The paper must state an identification, structural, trial, or sensitivity contract that explains how alternative outcomes are learned or assessed \citep{Pearl2009Causality,PearlMackenzie2018BookOfWhy}. This requirement makes Level~4 more restrictive than treatment-conditioned generation.

The lower boundary is Level~3. Sampling rollouts under two treatment tokens can reveal model sensitivity, but it does not show that their difference represents the treatment effect for the same patient. Counterfactual image editing, alternative explanations, and factual outcome prediction also remain outside Level~4 unless they estimate an unobserved clinical outcome under an alternative intervention. The upper boundary is not planning by default. A counterfactual estimator can compare interventions without choosing one, while a planner can select actions through a biased simulator without identifying counterfactual outcomes. Levels~4 and~5 answer distinct questions.

This category draws heavily from causal inference rather than from papers branded as world models. That is analytically important. The methods most explicit about potential outcomes, time-varying confounding, and treatment support often do not use world-model language, while systems described as world models frequently emphasize generation or control. The gap between these lineages is not a reason to relax the category. It identifies a missing bridge between flexible patient simulation and credible alternative-outcome estimation.

\subsection{Clinical Evidence Landscape}

The most natural Level~4 applications are longitudinal treatment-effect estimation, survival analysis, critical care, and oncology. These settings compare treatment sequences, doses, or policies over time and must address treatment-dependent censoring, competing risks, missingness, and evolving confounders. Representative methods include counterfactual recurrent networks, causal transformers, G-computation networks, and time-series deconfounding approaches \citep{Bica2020CRN,Melnychuk2022CausalTransformer,Li2020GNet,Seedat2022TECDE}. Their principal object is an alternative outcome trajectory rather than a realistic sample from the observed care distribution.

The broader treatment-effect cluster includes doubly robust representation learning, disentangled recurrent networks, heterogeneous-effect ODE discovery, and longitudinal matching through synthetic twins \citep{Zeng2020DoubleRobustRepresentation,Berrevoets2021DisentangledCounterfactualRecurrent,Kacprzyk2024OdeDiscoveryLongitudinal,Qian2021SynctwinTreatmentEffect}. Systems-based pharmacology twins add a mechanistic route to dose--response analysis \citep{Susilo2023Systemsbaseddigital}. These methods vary in representation and assumptions, but they share an explicit comparison target that is absent from descriptive action-conditioned generation.

Survival models provide a particularly clear estimand. Treatment-specific survival functions or potential-outcome regressors can compare clinically defined alternatives while accounting for censoring and treatment assignment. Open packages such as auton-survival illustrate how counterfactual estimation can be made inspectable, though software availability does not by itself guarantee identification. Oncology is a prominent application because treatment decisions and delayed outcomes are central, but the same framework applies to critical care, chronic disease, and other longitudinal interventions.

Time-varying clinical applications make the assumptions harder to avoid. Counterfactual sepsis-treatment frameworks, recurrent marginal structural networks, time-series deconfounding, G-transformers, and balanced survival models address different combinations of dynamic treatment, hidden or measured confounding, and censoring \citep{Jeter2021LearningTreatHypotensive,Lim2018ForecastingTreatmentResponses,Bica2020TimeSeriesDeconfounder,Xiong2024GTransformerCounterfactual,Chapfuwa2021EnablingCounterfactualSurvival,Nagpal2022autonsurvivalOpen}. Their outputs should be read at the level of the stated estimand and support assumptions, not as unrestricted patient-specific futures.

Medical imaging reaches Level~4 less often. A model that generates two plausible post-treatment images still faces the missing-outcome problem: only one trajectory is observed for a patient. Images can provide rich intermediate states, but they do not remove confounding or establish an outcome contrast. Mechanistic tumor or physiological models may offer structural support when their parameters and intervention pathways are validated. Their counterfactual interpretation remains bounded by model adequacy and patient-specific calibration.

Two examples show the range of possible anchors. Latent causal modeling can construct counterfactual brain MRI under an explicit structural account, while a digital-twin-guided virtual amiodarone test ties simulated intervention response to a clinically evaluated electrophysiology setting \citep{Peng2024LatentCausalModeling,Hwang2024ClinicalUsefulnessDigital}. Neither route removes the missing-potential-outcome problem. It instead makes the source of counterfactual support more visible: causal structure in one case and a calibrated mechanistic intervention pathway in the other.

The clinical evidence landscape is consequently narrower than the broader treatment-simulation literature. This scarcity should be reported as a finding rather than hidden by expansive terminology. It suggests that counterfactual capability currently resides mostly in an adjacent methodological tradition and that integration with multimodal world models remains incomplete.

This separation also changes how representative work should be selected. A treatment-conditioned generator is most informative here when it supplies an identified comparison, not simply because its outputs are visually rich. Conversely, a causal estimator belongs in the discussion even when it lacks world-model branding, because it directly addresses the alternative-outcome question that the level is meant to capture. The category is defined by the scientific estimand rather than the vocabulary of a research community.

\subsection{Representative Model Designs}

Balanced recurrent representations seek to remove information predictive of treatment assignment while preserving outcome-relevant history. Causal transformers extend this idea to long-range dependencies and time-varying treatments \citep{Bica2020CRN,Melnychuk2022CausalTransformer}. G-computation and related sequence models estimate outcomes by iterating conditional models under a specified intervention regime \citep{Li2020GNet}. Potential-outcome heads model treatment-specific outcomes directly, sometimes with propensity or weighting components. These designs make the estimand more explicit than a generic action-conditioned generator.

Graphical and structural causal models offer another route. They encode assumptions about how patient variables, treatments, measurements, and outcomes relate. Mechanistic simulators can play a similar role when biological equations specify how interventions alter dynamics. Neither structure is self-validating. Graphs require defensible causal assumptions; mechanistic systems require parameter identification and evidence that omitted processes do not dominate the decision. Hybrid models may combine flexible representation with a structured outcome or transition component, but the source of identification should remain visible.

Architecture is therefore secondary to the causal contract. A transformer described as causal does not qualify because of its name, and a small regression model can qualify when the estimand, assumptions, and validation are appropriate. The relevant design questions are which interventions are compared, over what horizon, for which population, under what support and consistency assumptions, and how time-varying confounding and censoring are handled.

\subsection{Evidence Quality and Failure Modes}

The defining validation problem is that both potential outcomes are never jointly observed for the same patient. Factual hold-out accuracy is insufficient because a model can predict observed outcomes well while estimating treatment contrasts poorly. Evidence may come from randomized or quasi-experimental variation, semi-synthetic data with known effects, trial emulation, mechanistic constraints, negative controls, held-out policy regimes, or sensitivity analyses. Each source supports a bounded claim; none turns an observational model into a universally valid patient-specific oracle.

Overlap and positivity are central. If a treatment was rarely given to patients in a particular state, the model must extrapolate to evaluate that alternative. Aggregate performance can hide this problem. Studies should report treatment support, effective sample sizes, uncertainty, and performance in weak-overlap regions. Time-varying settings also require attention to treatment-dependent censoring and post-treatment covariates. A model that conditions incorrectly on those variables can introduce bias even when its predictive metrics improve.

The common overclaim is to use the word counterfactual for any alternative generation. A clinically useful counterfactual statement must specify the estimand, intervention regime, target population, horizon, assumptions, and uncertainty. Benchmark counterfactual accuracy is valuable when the ground truth is known, but it does not automatically validate a real-patient claim. External data with different practice patterns, trial comparisons, and domain-specific sensitivity analysis can narrow that gap.

Another failure is excessive patient specificity. Population-average or subgroup treatment effects can be identified under assumptions that are insufficient for individual treatment effects. A model may output a patient-level number because its input is patient-specific, but the uncertainty in the individual contrast can remain large. Claims should match the level at which the evidence is informative.

\subsection{Level Takeaway and Next Threshold}

Level~4 is not simply a more elaborate generator. It changes the scientific target from an action-labeled future to an alternative outcome under an explicit comparison. The strongest evidence currently comes from causal sequence modeling, survival analysis, and mechanistic or semi-synthetic validation. These methods provide concepts that world-model systems need, while world models offer richer state and transition representations that causal estimators often lack.

\levelsummary
{Causal sequence, survival, structural, and mechanistic approaches provide explicit tools for estimating outcomes under alternative interventions, especially for longitudinal treatment questions.}
{Strict counterfactual evidence remains concentrated in the causal-inference lineage and is still weakly integrated with systems branded as medical world models. Generating alternatives is more common than identifying them.}
{The next threshold is not automatically Level~5. A decision system must operationally use a validated model to rank or select actions and must evaluate how transition, counterfactual, and uncertainty errors affect that choice.}
